\documentclass[11pt]{article}
\usepackage[margin=1in]{geometry}
\usepackage{amsmath,amssymb,amsthm,mathtools}
\usepackage{graphicx}
\usepackage{float}
\usepackage{microtype}
\usepackage[hidelinks]{hyperref}

\newtheorem{theorem}{Theorem}
\newtheorem{proposition}[theorem]{Proposition}
\newtheorem{corollary}[theorem]{Corollary}
\newcommand{\R}{\mathbb R}
\newcommand{\N}{\mathbb N}

\title{Length Rigidity and the Nowhere-Continuous-Inverse Obstruction\\
for Scottish Book Problem 155}
\author{Partial-result packet for arXiv:2308.03339}
\date{12 August 2026}

\begin{document}
\maketitle

\begin{center}
\fbox{\begin{minipage}{0.9\textwidth}
\textbf{Status: substantial partial result, likely valid; full problem open.}
A locally distance-preserving surjection between Banach spaces is globally
nonexpansive and preserves path length.  Continuity of the inverse at one
point forces a global affine isometry.  Thus every counterexample must have
a nowhere-continuous inverse and closed nowhere-dense local sheets.  An exact
locally isometric but globally contracting self-embedding of $\ell_\infty$
shows why surjectivity is the remaining essential obstacle.
\end{minipage}}
\end{center}

\section{Source question}

Mori \cite{Mori} studies Scottish Book Problem 155.  In modern notation it
asks: if $U:X\to Y$ is a bijection between real Banach spaces and for every
$x\in X$ there is $r_x>0$ such that $U|_{B(x,r_x)}$ preserves distances,
must $U$ preserve all distances?  The source proves this when $X$ is
separable and explicitly leaves the nonseparable case open.

\begin{figure}[H]
\centering
\includegraphics[width=0.92\textwidth]{figures/open_problem_crop.png}
\caption{The source question and its nonseparable context (PDF page 1).}
\end{figure}

Basso \cite{Basso} proves the arbitrary-density case under the stronger
assumption that every local restriction maps onto an \emph{open} target
neighborhood.  That hypothesis makes $U$ a covering/local-isometry map; it is
not automatic in the question above.

\section{Proof intuition}

The straight segment from $x$ to $y$ is compact, so finitely many of the
isometric neighborhoods control its image.  Its image has the same path
length, which immediately yields the one-sided global estimate
$\|U(x)-U(y)\|\leq\|x-y\|$.  The missing inequality is precisely a lifting
problem for target paths.

One open local image is enough to remove that obstruction: Mankiewicz's
extension theorem supplies a global affine isometry, and agreement propagates
through connectedness.  Hence one continuity point of $U^{-1}$ settles the
whole problem.  If there is no such point, all complete local images must be
nowhere dense.  The final construction shows that thin local images can
indeed fold distant points close together, even in $\ell_\infty$; it fails
only the source's surjectivity requirement.

\section{Length rigidity and a one-point criterion}

The next statement allows surjectivity without injectivity until the inverse
is mentioned.

\begin{theorem}\label{thm:structure}
Let $X,Y$ be real Banach spaces and let $U:X\to Y$ be a surjection such that
every $x\in X$ has a radius $r_x>0$ for which $U|_{B(x,r_x)}$ is an
isometric embedding.  Then:
\begin{enumerate}
\item $U$ preserves the length of every rectifiable path and is globally
      $1$-Lipschitz;
\item if some local image $U(B(x,r))$, on which $U$ is isometric, has
      nonempty interior, then $U$ is a global affine surjective isometry;
\item if $U$ is bijective and $U^{-1}$ is continuous at one point of $Y$,
      then $U$ is a global affine surjective isometry.
\end{enumerate}
\end{theorem}

\begin{proof}
Let $\gamma:[0,1]\to X$ be a rectifiable path.  Compactness of its image and
the local hypothesis give a subdivision fine enough that each subarc lies in
one isometric ball.  Refining any partition in this way shows that the sums
of successive distances used to define the lengths of $\gamma$ and
$U\circ\gamma$ agree.  Taking suprema gives
$L(U\circ\gamma)=L(\gamma)$.  Applying this to the linear segment from $x$
to $y$ yields
\[
 \|U(x)-U(y)\|\leq L(U\circ\gamma)=L(\gamma)=\|x-y\|.
\]

For (2), Mori's Lemma 2, which is an application of Mankiewicz's extension
theorem, says that the local restriction extends uniquely to an affine
surjective isometry $A:X\to Y$.  Let $D$ be the set of points near which
$U=A$.  It is nonempty and open.  If $D\ne X$, choose a boundary point $z$.
On a sufficiently small isometric ball at $z$, the maps agree on a nonempty
open subset coming from $D$.  Consequently the image of that local ball has
nonempty interior; its unique affine extension must be $A$.  Hence $U=A$ in
a neighborhood of $z$, a contradiction.  Connectedness of $X$ gives
$D=X$.

For (3), write $y=U(x)$ for a continuity point of $U^{-1}$ and choose a
local isometric ball $B(x,r)$.  Continuity at $y$ provides $\delta>0$ with
\[
 U^{-1}(B(y,\delta))\subset B(x,r).
\]
Surjectivity and injectivity therefore give
$B(y,\delta)\subset U(B(x,r))$.  Apply (2).
\end{proof}

\begin{corollary}\label{cor:b1}
Under the hypotheses of Theorem \ref{thm:structure}, if $U$ is bijective and
$U^{-1}:Y\to X$ is of Baire class one, then $U$ is a global affine
surjective isometry.
\end{corollary}

\begin{proof}
A Baire-one map from the completely metrizable space $Y$ to a metric space
has a point of continuity.  Use Theorem \ref{thm:structure}(3).
\end{proof}

\section{The category obstruction}

Let $\operatorname{cov}_{\mathrm{nd}}(Y)$ be the least cardinality of a cover
of $Y$ by closed nowhere-dense sets.  Let $\operatorname{dens}(X)$ denote
the density character of $X$.

\begin{proposition}\label{prop:category}
If
\[
 \operatorname{dens}(X)<\operatorname{cov}_{\mathrm{nd}}(Y),
\]
then every map in Theorem \ref{thm:structure} is a global affine surjective
isometry.
\end{proposition}

\begin{proof}
For each $x$, first replace its local radius $r_x$ by the open ball
$B^\circ(x,r_x/4)$.  These smaller balls cover $X$.  A metric space of
density $\kappa$ has Lindel\"of number at most $\kappa$, so select a subcover
of cardinal at most $\operatorname{dens}(X)$.  Replace each selected
$B^\circ(x,r_x/4)$ by $B(x,r_x/2)$; the closed balls still cover $X$ and lie
inside isometric neighborhoods.  The image of each is closed because it is
isometric to a complete space.  If all these images were nowhere dense, their
union could not cover $Y$ by the displayed cardinal inequality.  Therefore
one has interior, and Theorem \ref{thm:structure}(2) applies.
\end{proof}

For separable $X$, Proposition \ref{prop:category} reduces to the ordinary
Baire-category argument in the source.  Its point here is structural: in any
genuine nonseparable counterexample, every complete local sheet is closed
nowhere dense, and enough such sheets must cover $Y$.

\section{A sharp near-counterexample on the bounded-sequence space}

Define the one-periodic tent wave
\[
 \tau(t)=\operatorname{dist}(t,\mathbb Z)
\]
and
\[
 F(t)=\bigl(\tau(t),\tau(t+1/4),\tfrac1{10}\arctan t\bigr)
 \in\ell_\infty^3.
\]

\begin{proposition}\label{prop:tent}
There is an injective self-map $T:\ell_\infty\to\ell_\infty$ whose
restriction to every ball of radius $1/16$ is isometric, but $T$ is not a
global isometric embedding.  In particular, the surjectivity hypothesis in
Scottish Book Problem 155 cannot be discarded.
\end{proposition}

\begin{proof}
The corner sets of $\tau(t)$ and $\tau(t+1/4)$ are respectively
$\frac12\mathbb Z$ and $\frac12\mathbb Z-\frac14$.  For every $a\in\R$,
one of these two sets is at distance at least $1/8$ from $a$.  The
corresponding tent coordinate is affine with slope $1$ or $-1$ throughout
the radius-$1/16$ neighborhood of $a$.  All three coordinates of $F$ are
$1$-Lipschitz, while that chosen coordinate witnesses equality.  Thus
\[
 \|F(s)-F(t)\|_\infty=|s-t|
 \quad(s,t\in B(a,1/16)).                              \tag{1}
\]
The last coordinate makes $F$ injective.  On the other hand, for integers
$N>0$,
\[
 \|F(N)-F(0)\|_\infty=\tfrac1{10}\arctan N<N=|N-0|.   \tag{2}
\]

Apply $F$ coordinatewise:
\[
 \mathcal F((x_i)_{i\in\N})=(F_j(x_i))_{(i,j)\in\N\times\{1,2,3\}}.
\]
Equation (1), uniformly in $i$, shows that $\mathcal F$ is isometric on
every radius-$1/16$ ball of $\ell_\infty$; it is injective, and (2) applied
to constant sequences shows global contraction.  Finally reindex
$\N\times\{1,2,3\}$ by $\N$ to obtain the claimed self-map $T$.
\end{proof}

The range of $T$ satisfies a coordinatewise curve constraint and is not all
of $\ell_\infty$.  Attempts to absorb this constraint by shifts or triangular
shears either preserved a non-surjective range condition or destroyed local
sup-distance equality.  This is the remaining gap between the construction
and a counterexample to the source problem.

\section{Verification, scope, and novelty}

The accompanying script checked 20,000 pseudorandom pairs in common
radius-$1/16$ neighborhoods and the global contraction in (2).  This is only
a transcription audit; Proposition \ref{prop:tent} is proved exactly above.

Eight focused upgrade attempts covered length lifting, convex bicombings,
midpoint characterizations, inverse continuity, generalized category,
Baire-one regularity, the tent-wave construction, and surjective absorption.
The full bijective nonseparable problem remains open because no argument
forces a target segment to lift and no locally isometric thin embedding was
made surjective.

On 12 August 2026, bounded searches of the four run indexes and current
literature used the source id, Problem 155, Mazur--Sternbach, the
nonseparable case, and later citations.  They found no full resolution.
Basso's result uses the strictly stronger open-image hypothesis.
Novelty confidence for the elementary partial criteria is moderate-low;
validity confidence is high.

\paragraph{Human-review recommendation.}
Review as a partial result.  The most useful checks are the closed-ball
subcover sentence in Proposition \ref{prop:category}, the one-point inverse
continuity reduction, and the uniform radius in Proposition
\ref{prop:tent}.

\begin{thebibliography}{9}
\bibitem{Mori}
M.~Mori,
\emph{On the Scottish Book Problem 155 by Mazur and Sternbach},
arXiv:2308.03339, 2023; C. R. Math. 362 (2024), 813--816.

\bibitem{Basso}
G.~Basso,
\emph{On a problem of Mazur and Sternbach},
arXiv:2308.08400v2, 2023.
\end{thebibliography}

\end{document}
